% =====================================================================
% TA-IND-01 — Informe Académico Individual (Aplicaciones Distribuidas)
% Análisis crítico de ClickHouse como sistema distribuido y su
% comparación con Elasticsearch
%
% Compilar desde la carpeta documento/, en este orden:
%   pdflatex main.tex
%   biber main
%   pdflatex main.tex
%   pdflatex main.tex
% =====================================================================
\documentclass[12pt,a4paper]{article}

\usepackage[T1]{fontenc}
\usepackage[utf8]{inputenc}
\usepackage[spanish,es-tabla,es-nodecimaldot]{babel}
\usepackage{lmodern}
\usepackage[margin=2.5cm]{geometry}
\usepackage{graphicx}
\graphicspath{{figuras/}{../capturas/}{../anexos/}}
\usepackage{booktabs}
\usepackage{tabularx}
\usepackage{float}
\usepackage{caption}
\usepackage{xcolor}
\usepackage{listings}
\lstset{basicstyle=\ttfamily\small, breaklines=true, frame=single, columns=fullflexible}
\usepackage{csquotes}
\usepackage[backend=biber,style=ieee]{biblatex}
\addbibresource{referencias.bib}
\usepackage[titletoc,title]{appendix}
\addto\captionsspanish{%
  \renewcommand{\appendixname}{Anexo}%
  \renewcommand{\appendixtocname}{Anexos}%
  \renewcommand{\appendixpagename}{Anexos}%
}
\usepackage[hidelinks]{hyperref}
\usepackage{xurl}

\begin{document}

% ---------------------------------------------------------------------
% CARÁTULA (verifica facultad, nivel y docente antes de entregar)
% ---------------------------------------------------------------------
\begin{titlepage}
  \centering
  {\large\bfseries UNIVERSIDAD TÉCNICA ESTATAL DE QUEVEDO\par}
  \vspace{0.3cm}
  {\large Facultad de Ciencias de la Computación y Diseño Digital\par}
  \vspace{0.2cm}
  {\large Carrera de Software\par}
  \vspace{1.5cm}
  {\Large\bfseries Aplicaciones Distribuidas\par}
  \vspace{0.4cm}
  {\large TA-IND-01 — Informe Académico Individual\par}
  \vspace{1.5cm}
  {\LARGE\bfseries Análisis crítico de ClickHouse como sistema distribuido y su comparación con Elasticsearch\par}
  \vfill
  \begin{tabular}{rl}
    \textbf{Estudiante:} & Washington German Apunte Pazmiño \\
    \textbf{Nivel y paralelo:} & Séptimo nivel, paralelo B \\
    \textbf{Docente:} & Gleiston Guerrero Ulloa, PhD \\
    \textbf{Fecha:} & 15 de noviembre de 2026 \\
  \end{tabular}
  \par\vspace{1cm}
  \textbf{Repositorio:}\par
  \url{https://github.com/Washinton-Apunte/TA-IND-01-ClickHouse}
\end{titlepage}

% ---------------------------------------------------------------------
% RESUMEN (150 a 200 palabras). Escríbelo al final, con resultados.
% ---------------------------------------------------------------------
\pagenumbering{roman}
\begin{abstract}
% Sistema analizado, enfoque y conclusiones principales.
\end{abstract}

\tableofcontents
\newpage
\pagenumbering{arabic}

% Mínimo 5 páginas de contenido (sin carátula, anexos ni referencias).

\section{Introducción}
% Qué problema resuelve ClickHouse, quién lo usa, a qué escala opera
% y por qué lo elegiste.

\section{Arquitectura del sistema}
% Componentes y cómo se comunican. Incluye tu diagrama (draw.io o TikZ,
% hecho por ti) y cómo evolucionó desde el bosquejo a mano (Anexo A).

\section{Análisis con los conceptos de la Unidad 1}

\subsection{Transparencias de distribución}
% Acceso, ubicación, migración, reubicación, replicación, concurrencia
% y fallos: cuáles ofrece ClickHouse, cuáles no y por qué.

\subsection{CAP y PACELC}
% Qué prioriza ante una partición de red y qué prioriza sin partición
% (latencia frente a consistencia), con evidencia de la documentación
% técnica o de artículos.

\subsection{Modelos de comunicación}
% Cliente-servidor, peer-to-peer o híbrido; sincrónica o asincrónica;
% sockets, RPC o mensajería.

\subsection{Tolerancia a fallos}
% Qué fallos soporta, cómo los detecta y cómo se recupera (replicación,
% consenso, reintentos), con al menos un incidente real documentado.

\section{Experimentación propia}
% Capturas de la ejecución (evidencia 1) y experimento (evidencia 2),
% con su interpretación a la luz de la sección anterior.

\section{Ejemplos del entorno del estudiante}
% Al menos dos sistemas reales que usas o conoces de primera mano,
% relacionados con un concepto de la unidad (evidencia 3).

\section{Análisis comparativo}
% Tabla ClickHouse frente a Elasticsearch con al menos cinco criterios
% de la Unidad 1, y su discusión.

\section{Postura crítica}
% Fortalezas, limitaciones, compromisos (trade-offs) y contextos en los
% que ClickHouse no sería la mejor opción. Todo argumentado y respaldado.

\section{Conclusiones}
% Afirmaciones concretas que salen del análisis y del experimento.

\section*{Declaración de uso de inteligencia artificial}
\addcontentsline{toc}{section}{Declaración de uso de inteligencia artificial}
% Herramienta, versión, para qué se usó y cómo verificaste lo obtenido
% (o declaración de que no se usó).

\printbibliography[title={Referencias},heading=bibintoc]

\begin{appendices}
\section{Esquema inicial a mano}\label{anexo:a}
% Fotografía del primer bosquejo con tu nombre y la fecha visibles.

\section{Capturas de lectura de las fuentes}\label{anexo:b}
% Al menos 3 capturas con fragmentos subrayados o anotados y la página;
% cada una debe corresponder a una cita del texto.
\end{appendices}

\end{document}
